\clearpage
\fancyhead[L]{\small BEOBACHTUNGEN.}
\hypertarget{native-466-letter}{}
\begin{center}
\rule{0.55\linewidth}{0.4pt}\\[2pt]
Astronomische Nachrichten Band VII. Nr. 145. Seite 15. 1828 September.\\[2pt]
\rule{0.55\linewidth}{0.4pt}
\end{center}

\begin{center}[\textsc{Gauss} an \textsc{Schumacher.}]\quad 1827 Oct. 11.\end{center}

Hier meine bisherigen \emph{Sonnenzenithdistanzen}, die nur von der Refraction befreit sind. In der folgenden Columne steht die Parallaxe (sie ist für das Centrum berechnet, oder vielmehr aus der ausgefüllten Columne Hülfstafeln 1821 entlehnt; es ist genauer sie für den beobachteten Rand zu nehmen, der Unterschied ist aber nur $\pm0''03$ und gleicht sich aus), dann die mit der Polhöhe $51^\circ31'48''$, und dem Sonnenhalbmesser aus den Hülfstafeln von 1821 berechnete Declination, dann die Differenz mit den Hülfstafeln, dann die Differenz wenn, wie aus allen sich ergibt, $\mathrm{d}r=-0''77$ gesetzt wird.

\begin{center}
\begingroup\fontsize{10}{12}\selectfont\setlength{\tabcolsep}{3pt}
\begin{tabular}{l|c|l|l|l|l|l}
1827 Sept. 11 & U & $47^\circ0'51''03$ & $6''21$ & $+4^\circ46'58''75$ & $+3''65-\mathrm{d}r$ & $+4''42$ \\
\hfill 15 & U & $48\;32\;42.39$ & $6.37$ & $+3\;15\;8.57$ & $+1.85-\mathrm{d}r$ & $+2.60$ \\
\hfill 21 & O & $50\;20\;23.69$ & $6.61$ & $+0\;55\;32.74$ & $+5.36+\mathrm{d}r$ & $+4.59$ \\
\hfill 22 & O & $50\;43\;46.00$ & $6.65$ & $+0\;32\;10.20$ & $+4.90+\mathrm{d}r$ & $+4.13$ \\
\hfill 23 & O & $51\;7\;8.48$ & $6.70$ & $+0\;8\;47.50$ & $+3.40+\mathrm{d}r$ & $+2.67$ \\
\hfill Oct. 3 & U & $55\;33\;0.97$ & $7.05$ & $-3\;45\;4.43$ & $+3.63-\mathrm{d}r$ & $+4.40$ \\
\hfill 4 & U & $55\;56\;13.95$ & $7.08$ & $-4\;8\;17.11$ & $+1.91-\mathrm{d}r$ & $+2.68$ \\
\hfill 5 & O & $55\;47\;23.33$ & $7.12$ & $-4\;31\;30.25$ & $+3.55+\mathrm{d}r$ & $+2.78$ \\
\end{tabular}
\endgroup
\end{center}

Die $0''77$, um welche ich den Sonnenhalbmesser vermindern muss, können vielleicht zum Theile davon abhängen, dass bei dem Urtheil über die Halbirung des Fadenintervalls, wo die eine Hälfte hell, die andere dunkel ist, etwas individuelles mit unterlaufen mag.

Ich übersende Ihnen ferner hier meine Ceresbeobachtungen.

\begin{center}
\begin{tabular}{l|l|l|l}
1827 Sept. 27 & $12^{\rm u}11^{\rm m}28^{\rm s}7$ M. Z. & $8^\circ50'36''1$ & $-13^\circ14'56''1$ \\
\hfill Oct. 3 & $11\;42\;50.2$ & $7\;34\;37.6$ & $13\;41\;34.8$ \\
\hfill 4 & $11\;38\;3.7$ & $7\;21\;56.4$ & $13\;45\;6.0$ \\
\hfill 5 & $11\;33\;17.7$ & $7\;9\;22.8$ & $-13\;48\;24.9$ \\
\hfill 6 & $11\;28\;31.9$ & $9\;56\;52.9$ & $-13\;51\;37.1$ \\
\end{tabular}
\end{center}

\clearpage
\fancyhead[L]{\small CERES PALLAS.}
\SourceLeaf{467}
\begin{center}
\rule{0.55\linewidth}{0.4pt}\\[2pt]
Astronomische Nachrichten Band VII. Nr. 161. Seite 329. 1829 Juni.\\[2pt]
\rule{0.55\linewidth}{0.4pt}
\end{center}

\begin{center}[Beobachtungen der \emph{Ceres} in Göttingen von \textsc{Gauss.}]\end{center}

\begin{center}
\begin{tabular}{l|l|l|l}
1829 & Mittlere Zeit & Gerade Aufsteig. & Abweich. Nordl. \\
\hline
Januar 22 & $10^{\rm u}21^{\rm m}46^{\rm s}2$ & $97^\circ24'25''0$ & $30^\circ22'0''4$ \\
\hfill 25 & $10\;7\;41.4$ & $96\;50\;4.0$ & $30\;30\;8.3$ \\
Februar 11 & $8\;52\;51.7$ & $94\;49\;55.8$ & $30\;58\;55.4$ \\
\hfill 19 & $8\;20\;52.8$ & $94\;41\;58.8$ & $31\;4\;46.7$ \\
März 5 & $7\;29\;47.9$ & $95\;41\;35.9$ & $31\;6\;19.2$ \\
\hfill 14 & $6\;59\;57.5$ & $97\;5\;4.0$ & $31\;2\;42.0$ \\
\hfill 16 & $6\;53\;37.2$ & $97\;27\;56.7$ & \\
\hfill 18 & $6\;47\;22.3$ & $97\;52\;14.5$ & $30\;59\;55.9$ \\
\end{tabular}
\end{center}

Die Beobachtung der Opposition war durch die anhaltend ungünstige Witterung vereitelt.

\clearpage
\fancyhead[L]{\small CERES PALLAS.}
\hypertarget{native-467-letter}{}
\begin{center}
\rule{0.55\linewidth}{0.4pt}\\[2pt]
Astronomische Nachrichten Band VIII. Nr. 185. Seite 321. 1830 August.\\[2pt]
\rule{0.55\linewidth}{0.4pt}
\end{center}

\begin{center}Auszug aus einem Schreiben des Herrn Hofraths \textsc{Gauss} an den Herausgeber.\quad Göttingen 1830. Juli 20.\end{center}

Hiebei übersende ich Ihnen meine bei den diesjährigen Oppositionen angestellten Beobachtungen der \emph{Ceres} und Pallas.

\begin{center}
\emph{Ceres.}

\begin{tabular}{l|l|l|l}
1830 & Mittlere Zeit & Gerade Aufsteig. & Abweichung \\
\hline
April 24 & $12^{\rm u}39^{\rm m}21^{\rm s}9$ & $222^\circ20'25''8$ & $4^\circ18'27''6$ S. \\
\hfill 25 & $34\;33.0$ & $7\;9.1$ & $16\;45.0$ \\
\hfill 26 & $29\;44.3$ & $221\;53\;54.1$ & $15\;9.0$ \\
\hfill 27 & $24\;54.7$ & $40\;27.3$ & $13\;38.8$ \\
\hfill 28 & $20\;5.0$ & $26\;58.3$ & $12\;14.6$ \\
\hfill 29 & $15\;15.3$ & $13\;28.8$ & $10\;52.4$ \\
\hfill 30 & $10\;25.1$ & $220\;59\;52.5$ & $9\;37.2$ \\
Mai 2 & $0\;44.7$ & $32\;39.1$ & $7\;29.2$ \\
\hfill 5 & $11\;46\;14.9$ & $219\;52\;0.7$ & $5\;2.0$ \\
\end{tabular}
\end{center}

\begin{center}
Beobachtungen der \emph{Pallas} in Göttingen.

\begin{tabular}{l|l|l|l}
1830 & Mittlere Zeit & Gerade Aufsteig. & Abweichung \\
\hline
April 24 & $13^{\rm u}0^{\rm m}9^{\rm s}1$ & $227^\circ33'5''2$ & $22^\circ38'1''6$ N. \\
\hfill 25 & $12\;55\;26.1$ & $21\;17.4$ & $49\;48.7$ \\
\hfill 26 & $50\;42.4$ & $9\;17.7$ & $23\;1\;16.3$ \\
\hfill 27 & $45\;58.5$ & $226\;57\;15.4$ & $12\;19.4$ \\
\hfill 28 & $41\;14.1$ & $45\;6.1$ & $22\;58.4$ \\
\hfill 29 & $36\;29.3$ & $32\;51.7$ & $33\;23.1$ \\
\hfill 30 & $81\;44.3$ & $20\;33.1$ & $43\;21.8$ \\
\end{tabular}
\end{center}

Von der neulich bei Tage gesehenen Bedeckung des Aldebaran habe ich den Eintritt sehr gut beobachtet:

\begin{center}Juli 16. Eintritt $\alpha$ Tauri am hellen M. R. $0^{\rm u}44^{\rm m}44^{\rm s}0$ M. Z.\end{center}

Gleich nachher verdunkelten Wolken den Himmel, und der Mond kam nicht wieder zu Gesicht.
